% Propietario conservado: /Users/ruben/Documents/ChatGPT/jueces y controles/REVISION_ARGUMENTAL_APERTURAS_CIERRES_20260915/II/manuscript_es/sections/07_elipse.tex
\section{Ellipse d’action et réciprocité modulaire}
\label{ii:sec:ellipse} Les sorties $A,C^*,\hbar>0$ étant fixées dans la carte $|C^*/A|<1$, définissons
\[
\eta_{\rm el}=\operatorname{artanh}(C^*/A),\quad
a_S=\sqrt{2\hbar}\,e^{\eta_{\rm el}/2},\quad
b_S=\sqrt{2\hbar}\,e^{-\eta_{\rm el}/2}.
\]
Les coordonnées canoniques équilibrées ont pour unité la racine de l’action ; on n’additionne pas ici position et quantité de mouvement sans carte dimensionnelle. On a
\[
a_Sb_S=2\hbar,\qquad
\tau_{\rm el}=i\,\frac{b_S}{a_S},\qquad
R_{\rm el}=\sqrt{\frac{b_S}{a_S}}
=\left(\frac{A-C^*}{A+C^*}\right)^{1/4}.
\]
La racine positive détermine la composition $\tau_{\rm el}=iR_{\rm el}^2$,
et permet de retrouver le rapport angulaire :
\[
\frac{C^*}{A}=\frac{1-R_{\rm el}^4}{1+R_{\rm el}^4}.
\]
L’inversion de forme envoie $R_{\rm el}$ sur $R_{\rm el}^{-1}$ et $\tau_{\rm el}$ sur $-1/\tau_{\rm el}$ ; l’équation~\eqref{ii:eq:ii-elipse-reciproca} développe la réciprocité de la forme.

\begin{proposition}[Énergie coïncidente, action discernable]
Soit
\[
D(\eta)=\begin{pmatrix}e^\eta&0\\0&e^{-\eta}\end{pmatrix},
\quad
S_2=\begin{pmatrix}0&1\\1&0\end{pmatrix},
\quad
J_2=\begin{pmatrix}0&-1\\1&0\end{pmatrix}.
\]
Les deux applications satisfont $T^{\mathsf T}D(-\eta)T=D(\eta)$.
Cependant, pour $\omega=\dd Q\wedge\dd P$,
\[
S_2^*\omega=-\omega,\qquad J_2^*\omega=\omega.
\]
Si $\mathcal I(\gamma)=\oint_\gamma P\,\dd Q$ sur une courbe fermée,
\[
\mathcal I(S_2\gamma)=-\mathcal I(\gamma),\qquad
\mathcal I(J_2\gamma)=\mathcal I(\gamma).
\]
\end{proposition}
\begin{proof}
La conjugaison échange les entrées diagonales de $D$.
Les déterminants de $S_2,J_2$ sont $-1,1$, respectivement, ce qui
prouve la transformation de $\omega$. Pour $\lambda=P\,\dd Q$,
\[
S_2^*\lambda=\dd(PQ)-\lambda,\qquad
J_2^*\lambda=\lambda-\dd(PQ).
\]
Intégrer la différentielle exacte sur la courbe fermée complète la preuve.
\end{proof}

L’énergie quadratique et le module ne suffisent pas à distinguer les deux
transports. L’orientation de l’action les distingue, même
dans la forme circulaire. Le marquage géométrique fait ainsi partie de la lecture
récupérable, ce n’est pas un ornement ajouté à la valeur.

